% =========================================================
% Proof: Closed Sets Contain Their Boundary
% Source: notes/open-closed-sets/open-closed-sets.tex
% Difficulty: Medium
% =========================================================

\subsection*{Closed Sets Contain Their Boundary}
\label{prf:closed-owns-boundary}

\begin{remark*}[Return]
$\leftarrow$ \hyperref[prop:closed-owns-boundary]{Back to Proposition in Notes}
\end{remark*}

\begin{proof}
\Claim $F \subseteq X$ is closed if and only if $\partial F \subseteq F$.

\Given A metric space $(X,d)$ and $F \subseteq X$.

\Goal To prove the biconditional.

\Strategy Use the decomposition $X = \operatorname{int}(F) \sqcup \partial F \sqcup \operatorname{ext}(F)$.
$F$ is closed iff $X \setminus F$ is open iff $X \setminus F = \operatorname{ext}(F)$
iff $\partial F \subseteq F$.
Make each step precise.

% -------------------------------------------------------
% YOUR PROOF HERE
% -------------------------------------------------------

\end{proof}

\begin{remark*}[Proof shape]
The decomposition theorem turns closedness into a statement about
which partition class the boundary points fall into.
\end{remark*}

\begin{remark*}[Dependencies]
\begin{itemize}
  \item Theorem~\ref{thm:ibe-decomp}: interior--boundary--exterior decomposition.
  \item Definition~\ref{def:closed-set}: $F$ closed iff $X \setminus F$ open.
  \item Definition~\ref{def:exterior}: $\operatorname{ext}(F) = \operatorname{int}(X \setminus F)$.
\end{itemize}
\end{remark*}
